%! Author = nursace
%! Date = 16.09.2026

\documentclass[aspectratio=169,11pt]{beamer}
\usepackage{color}
% Week 1 Lecture — What is AI / ML? Course setup
% Course: AI Frameworks (ages 16--18; intuition over math)

% Shared Beamer preamble for AI Frameworks
% Audience: college students ages 16--18 (intuition over math)
% Include with: % Shared Beamer preamble for AI Frameworks
% Audience: college students ages 16--18 (intuition over math)
% Include with: % Shared Beamer preamble for AI Frameworks
% Audience: college students ages 16--18 (intuition over math)
% Include with: \input{preamble}

\usetheme{Madrid}
\usecolortheme{default}

% Clean palette: deep slate + teal accent (distinct from Java navy decks)
\definecolor{LectureNavy}{HTML}{1A365D}
\definecolor{LectureAccent}{HTML}{319795}
\definecolor{LectureGray}{HTML}{4A5568}
\definecolor{CodeBg}{HTML}{F7FAFC}
\definecolor{AlertBg}{HTML}{FFF5F5}
\definecolor{TryBg}{HTML}{F0FFF4}
\definecolor{QuizBg}{HTML}{EBF8FF}
\definecolor{AnswerKeyBg}{HTML}{FFFFF0}
\definecolor{AnswerGold}{HTML}{B7791F}
\definecolor{KeywordBlue}{HTML}{2B6CB0}
\definecolor{StringGreen}{HTML}{276749}
\definecolor{CommentGray}{HTML}{718096}

\setbeamercolor{structure}{fg=LectureNavy}
\setbeamercolor{palette primary}{bg=LectureNavy,fg=white}
\setbeamercolor{palette secondary}{bg=LectureAccent,fg=white}
\setbeamercolor{palette tertiary}{bg=LectureGray,fg=white}
% Madrid title box is dark --- title/subtitle must be light for contrast
\setbeamercolor{title}{fg=white,bg=LectureNavy}
\setbeamercolor{subtitle}{fg=white,bg=LectureNavy}
\setbeamercolor{author}{fg=LectureNavy}
\setbeamercolor{date}{fg=LectureGray}
\setbeamercolor{institute}{fg=LectureGray}
\setbeamercolor{frametitle}{bg=LectureNavy,fg=white}
\setbeamercolor{block title}{bg=LectureNavy,fg=white}
\setbeamercolor{block body}{bg=CodeBg,fg=black}
\setbeamercolor{block title alerted}{bg=LectureAccent,fg=white}
\setbeamercolor{block body alerted}{bg=TryBg,fg=black}
\setbeamercolor{block title example}{bg=LectureGray,fg=white}
\setbeamercolor{block body example}{bg=AlertBg,fg=black}

\setbeamertemplate{navigation symbols}{}
\setbeamertemplate{itemize items}[circle]
\setbeamertemplate{enumerate items}[default]

\usepackage[T1]{fontenc}
\usepackage[utf8]{inputenc}
\usepackage{lmodern}
\usepackage{booktabs}
\usepackage{tabularx}
\usepackage{graphicx}
\usepackage{hyperref}
\usepackage{listings}
\usepackage{xcolor}
\usepackage{tikz}
\usetikzlibrary{positioning,arrows.meta,shapes.geometric}

% listings: Python without shell-escape
\lstdefinestyle{pythonstyle}{
    language=Python,
    basicstyle=\ttfamily\small,
    keywordstyle=\color{KeywordBlue}\bfseries,
    commentstyle=\color{CommentGray}\itshape,
    stringstyle=\color{StringGreen},
    numbers=left,
    numberstyle=\tiny\color{CommentGray},
    numbersep=6pt,
    frame=single,
    rulecolor=\color{LectureGray!40},
    backgroundcolor=\color{CodeBg},
    breaklines=true,
    showstringspaces=false,
    tabsize=4,
    captionpos=b,
}
\lstset{style=pythonstyle}

\newcommand{\pycode}[1]{{\ttfamily\small\detokenize{#1}}}
\newcommand{\trythis}[1]{%
    \begin{alertblock}{Try this}
    #1
    \end{alertblock}%
}
\newcommand{\commonmistake}[1]{%
    \begin{exampleblock}{Common mistake}
    #1
    \end{exampleblock}%
}
\newcommand{\remember}[1]{%
    \begin{block}{Remember}
    #1
    \end{block}%
}

% Formative in-class checks: question slide, then a dedicated Answer slide
\newcommand{\quizpause}{%
    \par\vspace{0.5em}
    {\small\textit{Pause --- answer (clicker / raise hand / neighbour) before the next slide.}}%
}
\newcommand{\quizbox}[1]{%
    \setbeamercolor{block title}{bg=KeywordBlue,fg=white}%
    \setbeamercolor{block body}{bg=QuizBg,fg=black}%
    \begin{block}{Quiz}
    #1
    \end{block}%
    \setbeamercolor{block title}{bg=LectureNavy,fg=white}%
    \setbeamercolor{block body}{bg=CodeBg,fg=black}%
}
\newcommand{\answerbox}[1]{%
    \setbeamercolor{block title}{bg=AnswerGold,fg=white}%
    \setbeamercolor{block body}{bg=AnswerKeyBg,fg=black}%
    \begin{block}{Answer}
    #1
    \end{block}%
    \setbeamercolor{block title}{bg=LectureNavy,fg=white}%
    \setbeamercolor{block body}{bg=CodeBg,fg=black}%
}

% Empty author/date placeholders for the lecturer to fill
\author{}
\date{}


\usetheme{Madrid}
\usecolortheme{default}

% Clean palette: deep slate + teal accent (distinct from Java navy decks)
\definecolor{LectureNavy}{HTML}{1A365D}
\definecolor{LectureAccent}{HTML}{319795}
\definecolor{LectureGray}{HTML}{4A5568}
\definecolor{CodeBg}{HTML}{F7FAFC}
\definecolor{AlertBg}{HTML}{FFF5F5}
\definecolor{TryBg}{HTML}{F0FFF4}
\definecolor{QuizBg}{HTML}{EBF8FF}
\definecolor{AnswerKeyBg}{HTML}{FFFFF0}
\definecolor{AnswerGold}{HTML}{B7791F}
\definecolor{KeywordBlue}{HTML}{2B6CB0}
\definecolor{StringGreen}{HTML}{276749}
\definecolor{CommentGray}{HTML}{718096}

\setbeamercolor{structure}{fg=LectureNavy}
\setbeamercolor{palette primary}{bg=LectureNavy,fg=white}
\setbeamercolor{palette secondary}{bg=LectureAccent,fg=white}
\setbeamercolor{palette tertiary}{bg=LectureGray,fg=white}
% Madrid title box is dark --- title/subtitle must be light for contrast
\setbeamercolor{title}{fg=white,bg=LectureNavy}
\setbeamercolor{subtitle}{fg=white,bg=LectureNavy}
\setbeamercolor{author}{fg=LectureNavy}
\setbeamercolor{date}{fg=LectureGray}
\setbeamercolor{institute}{fg=LectureGray}
\setbeamercolor{frametitle}{bg=LectureNavy,fg=white}
\setbeamercolor{block title}{bg=LectureNavy,fg=white}
\setbeamercolor{block body}{bg=CodeBg,fg=black}
\setbeamercolor{block title alerted}{bg=LectureAccent,fg=white}
\setbeamercolor{block body alerted}{bg=TryBg,fg=black}
\setbeamercolor{block title example}{bg=LectureGray,fg=white}
\setbeamercolor{block body example}{bg=AlertBg,fg=black}

\setbeamertemplate{navigation symbols}{}
\setbeamertemplate{itemize items}[circle]
\setbeamertemplate{enumerate items}[default]

\usepackage[T1]{fontenc}
\usepackage[utf8]{inputenc}
\usepackage{lmodern}
\usepackage{booktabs}
\usepackage{tabularx}
\usepackage{graphicx}
\usepackage{hyperref}
\usepackage{listings}
\usepackage{xcolor}
\usepackage{tikz}
\usetikzlibrary{positioning,arrows.meta,shapes.geometric}

% listings: Python without shell-escape
\lstdefinestyle{pythonstyle}{
    language=Python,
    basicstyle=\ttfamily\small,
    keywordstyle=\color{KeywordBlue}\bfseries,
    commentstyle=\color{CommentGray}\itshape,
    stringstyle=\color{StringGreen},
    numbers=left,
    numberstyle=\tiny\color{CommentGray},
    numbersep=6pt,
    frame=single,
    rulecolor=\color{LectureGray!40},
    backgroundcolor=\color{CodeBg},
    breaklines=true,
    showstringspaces=false,
    tabsize=4,
    captionpos=b,
}
\lstset{style=pythonstyle}

\newcommand{\pycode}[1]{{\ttfamily\small\detokenize{#1}}}
\newcommand{\trythis}[1]{%
    \begin{alertblock}{Try this}
    #1
    \end{alertblock}%
}
\newcommand{\commonmistake}[1]{%
    \begin{exampleblock}{Common mistake}
    #1
    \end{exampleblock}%
}
\newcommand{\remember}[1]{%
    \begin{block}{Remember}
    #1
    \end{block}%
}

% Formative in-class checks: question slide, then a dedicated Answer slide
\newcommand{\quizpause}{%
    \par\vspace{0.5em}
    {\small\textit{Pause --- answer (clicker / raise hand / neighbour) before the next slide.}}%
}
\newcommand{\quizbox}[1]{%
    \setbeamercolor{block title}{bg=KeywordBlue,fg=white}%
    \setbeamercolor{block body}{bg=QuizBg,fg=black}%
    \begin{block}{Quiz}
    #1
    \end{block}%
    \setbeamercolor{block title}{bg=LectureNavy,fg=white}%
    \setbeamercolor{block body}{bg=CodeBg,fg=black}%
}
\newcommand{\answerbox}[1]{%
    \setbeamercolor{block title}{bg=AnswerGold,fg=white}%
    \setbeamercolor{block body}{bg=AnswerKeyBg,fg=black}%
    \begin{block}{Answer}
    #1
    \end{block}%
    \setbeamercolor{block title}{bg=LectureNavy,fg=white}%
    \setbeamercolor{block body}{bg=CodeBg,fg=black}%
}

% Empty author/date placeholders for the lecturer to fill
\author{}
\date{}


\usetheme{Madrid}
\usecolortheme{default}

% Clean palette: deep slate + teal accent (distinct from Java navy decks)
\definecolor{LectureNavy}{HTML}{1A365D}
\definecolor{LectureAccent}{HTML}{319795}
\definecolor{LectureGray}{HTML}{4A5568}
\definecolor{CodeBg}{HTML}{F7FAFC}
\definecolor{AlertBg}{HTML}{FFF5F5}
\definecolor{TryBg}{HTML}{F0FFF4}
\definecolor{QuizBg}{HTML}{EBF8FF}
\definecolor{AnswerKeyBg}{HTML}{FFFFF0}
\definecolor{AnswerGold}{HTML}{B7791F}
\definecolor{KeywordBlue}{HTML}{2B6CB0}
\definecolor{StringGreen}{HTML}{276749}
\definecolor{CommentGray}{HTML}{718096}

\setbeamercolor{structure}{fg=LectureNavy}
\setbeamercolor{palette primary}{bg=LectureNavy,fg=white}
\setbeamercolor{palette secondary}{bg=LectureAccent,fg=white}
\setbeamercolor{palette tertiary}{bg=LectureGray,fg=white}
% Madrid title box is dark --- title/subtitle must be light for contrast
\setbeamercolor{title}{fg=white,bg=LectureNavy}
\setbeamercolor{subtitle}{fg=white,bg=LectureNavy}
\setbeamercolor{author}{fg=LectureNavy}
\setbeamercolor{date}{fg=LectureGray}
\setbeamercolor{institute}{fg=LectureGray}
\setbeamercolor{frametitle}{bg=LectureNavy,fg=white}
\setbeamercolor{block title}{bg=LectureNavy,fg=white}
\setbeamercolor{block body}{bg=CodeBg,fg=black}
\setbeamercolor{block title alerted}{bg=LectureAccent,fg=white}
\setbeamercolor{block body alerted}{bg=TryBg,fg=black}
\setbeamercolor{block title example}{bg=LectureGray,fg=white}
\setbeamercolor{block body example}{bg=AlertBg,fg=black}

\setbeamertemplate{navigation symbols}{}
\setbeamertemplate{itemize items}[circle]
\setbeamertemplate{enumerate items}[default]

\usepackage[T1]{fontenc}
\usepackage[utf8]{inputenc}
\usepackage{lmodern}
\usepackage{booktabs}
\usepackage{tabularx}
\usepackage{graphicx}
\usepackage{hyperref}
\usepackage{listings}
\usepackage{xcolor}
\usepackage{tikz}
\usetikzlibrary{positioning,arrows.meta,shapes.geometric}

% listings: Python without shell-escape
\lstdefinestyle{pythonstyle}{
    language=Python,
    basicstyle=\ttfamily\small,
    keywordstyle=\color{KeywordBlue}\bfseries,
    commentstyle=\color{CommentGray}\itshape,
    stringstyle=\color{StringGreen},
    numbers=left,
    numberstyle=\tiny\color{CommentGray},
    numbersep=6pt,
    frame=single,
    rulecolor=\color{LectureGray!40},
    backgroundcolor=\color{CodeBg},
    breaklines=true,
    showstringspaces=false,
    tabsize=4,
    captionpos=b,
}
\lstset{style=pythonstyle}

\newcommand{\pycode}[1]{{\ttfamily\small\detokenize{#1}}}
\newcommand{\trythis}[1]{%
    \begin{alertblock}{Try this}
    #1
    \end{alertblock}%
}
\newcommand{\commonmistake}[1]{%
    \begin{exampleblock}{Common mistake}
    #1
    \end{exampleblock}%
}
\newcommand{\remember}[1]{%
    \begin{block}{Remember}
    #1
    \end{block}%
}

% Formative in-class checks: question slide, then a dedicated Answer slide
\newcommand{\quizpause}{%
    \par\vspace{0.5em}
    {\small\textit{Pause --- answer (clicker / raise hand / neighbour) before the next slide.}}%
}
\newcommand{\quizbox}[1]{%
    \setbeamercolor{block title}{bg=KeywordBlue,fg=white}%
    \setbeamercolor{block body}{bg=QuizBg,fg=black}%
    \begin{block}{Quiz}
    #1
    \end{block}%
    \setbeamercolor{block title}{bg=LectureNavy,fg=white}%
    \setbeamercolor{block body}{bg=CodeBg,fg=black}%
}
\newcommand{\answerbox}[1]{%
    \setbeamercolor{block title}{bg=AnswerGold,fg=white}%
    \setbeamercolor{block body}{bg=AnswerKeyBg,fg=black}%
    \begin{block}{Answer}
    #1
    \end{block}%
    \setbeamercolor{block title}{bg=LectureNavy,fg=white}%
    \setbeamercolor{block body}{bg=CodeBg,fg=black}%
}

% Empty author/date placeholders for the lecturer to fill
\author{}
\date{}


\title{Week 1: What is AI / ML?}
\subtitle{AI Frameworks}

\begin{document}

%----------------------------------------------------------------------
    \begin{frame}
        \titlepage
    \end{frame}

%----------------------------------------------------------------------
    \begin{frame}{Today's learning goals}
        By the end of this lecture you should be able to:
        \begin{enumerate}
            \item Tell apart \textbf{AI}, \textbf{ML}, and \textbf{data science} in plain words
            \item Say what this course \emph{will} and \emph{won't} cover
            \item Explain \textbf{supervised learning} with examples, labels, and predictions
            \item Name our main tools: Python, Jupyter, and key libraries
        \end{enumerate}
    \end{frame}

    \begin{frame}{Agenda (80 minutes)}
        \begin{enumerate}
            \item Welcome and how the course works
            \item AI vs ML vs data science
            \item Scope: what we cover (and what we skip)
            \item Supervised learning intuition
            \item Tools overview: Python, Jupyter, libraries
            \item Looking ahead to practice
        \end{enumerate}
        \vspace{0.5em}
        Short quizzes are pauses to check understanding---not a test grade.
    \end{frame}

%======================================================================
    \section{Welcome}
%======================================================================

    \begin{frame}{Who this course is for}
        \begin{itemize}
            \item College students (about 16--18)
            \item Curious about AI, data, and building simple models
            \item Comfortable learning by running code and reading charts
        \end{itemize}
        \vspace{0.6em}
        \remember{We favour \textbf{intuition and practice} over heavy mathematics.
        You will see pictures and examples more often than formulas.}
    \end{frame}

    \begin{frame}{Weekly rhythm}
        \begin{columns}[T]
            \begin{column}{0.48\textwidth}
                \begin{block}{Lecture --- 80 min}
                    \begin{itemize}
                        \item Ideas and demos
                        \item Short quizzes
                        \item Questions welcome
                    \end{itemize}
                \end{block}
            \end{column}
            \begin{column}{0.48\textwidth}
                \begin{block}{Practice --- 80 min}
                    \begin{itemize}
                        \item Guided notebooks
                        \item Pair work
                        \item Done checklist
                    \end{itemize}
                \end{block}
            \end{column}
        \end{columns}
        \vspace{0.6em}
        Today after lecture: set up the environment and make your first predictions \emph{by hand}.
    \end{frame}

%======================================================================
    \section{AI vs ML vs data science}
%======================================================================

    \begin{frame}{Three overlapping ideas}
        Think of three circles that overlap---not three enemies.
        \vspace{0.8em}
        \begin{columns}[T]
            \begin{column}{0.32\textwidth}
                \textbf{AI}\\[0.3em]
                {\small Systems that seem to ``think'' or act intelligently (broad umbrella).}
            \end{column}
            \begin{column}{0.32\textwidth}
                \textbf{Machine learning}\\[0.3em]
                {\small Programs that \emph{learn patterns from data} instead of only hard-coded rules.}
            \end{column}
            \begin{column}{0.32\textwidth}
                \textbf{Data science}\\[0.3em]
                {\small Asking questions of data: clean, explore, visualise, sometimes model.}
            \end{column}
        \end{columns}
    \end{frame}

    \begin{frame}{A simple picture}
        \begin{center}
            \begin{tikzpicture}[
                every node/.style={align=center, font=\small},
                blob/.style={draw=LectureNavy, thick, rounded corners=6pt,
                minimum width=2.6cm, minimum height=1.5cm, fill=CodeBg}
            ]
                \node[blob, minimum width=9.2cm, minimum height=3.6cm, fill=LectureNavy!8]
                (ai) at (0,0) {};
                \node[font=\bfseries, LectureNavy] at (0,1.4) {Artificial Intelligence};
                \node[blob, minimum width=5.8cm, minimum height=2.2cm, fill=LectureAccent!15]
                (ml) at (0,-0.35) {};
                \node[font=\bfseries] at (0,0.35) {Machine Learning};
                \node[blob, fill=white] (ds) at (0,-0.55) {Data work\\\& modelling\\(this course)};
            \end{tikzpicture}
        \end{center}
        \vspace{0.3em}
        {\small Most of our semester lives inside the ML / data circle---practical modelling with tables.}
    \end{frame}

    \begin{frame}{Everyday examples (no math)}
        \begin{itemize}
            \item \textbf{Recommendations:} ``students who liked X also liked Y''
            \item \textbf{Spam filters:} learn from past spam vs not-spam emails
            \item \textbf{Photo tags:} the phone suggests who is in the picture
            \item \textbf{Weather apps:} models trained on historical measurements
        \end{itemize}
        \vspace{0.5em}
        Common pattern: \textbf{lots of examples} $\rightarrow$ \textbf{a system that generalises} to new cases.
    \end{frame}

    \begin{frame}{Rules vs learning}
        \begin{columns}[T]
            \begin{column}{0.48\textwidth}
                \begin{block}{Hand-written rules}
                    \begin{itemize}
                        \item \texttt{if} score $<$ 50 $\rightarrow$ fail
                        \item Brittle when life is messy
                        \item Easy to explain
                    \end{itemize}
                \end{block}
            \end{column}
            \begin{column}{0.48\textwidth}
                \begin{block}{Learning from data}
                    \begin{itemize}
                        \item Show many examples
                        \item Algorithm finds patterns
                        \item Needs careful checking
                    \end{itemize}
                \end{block}
            \end{column}
        \end{columns}
        \vspace{0.5em}
        This course teaches the second style---with tools that keep the workflow organised.
    \end{frame}

%--- Quiz 1 ---
    \begin{frame}{Quiz 1 --- AI / ML / data science}
        \quizbox{Which statement is the best match for \textbf{machine learning}?}
        \begin{enumerate}
            \item Designing web-page layout with HTML and CSS
            \item Writing only fixed \texttt{if--else} rules with no data
            \item Learning patterns from examples so we can make predictions on new cases
            \item Replacing all of mathematics with ChatGPT
        \end{enumerate}
        \quizpause
    \end{frame}

    \begin{frame}{Answer --- Quiz 1}
        \answerbox{\textbf{3.} Learning patterns from examples to predict on new cases.}
        \vspace{0.4em}
        AI is the broad umbrella; data science often includes exploring and communicating with data;
        ML is specifically about learning from data.
    \end{frame}

%======================================================================
    \section{Course scope}
%======================================================================

    \begin{frame}{What this course \emph{will} cover}
        \begin{itemize}
            \item Working with tables in \textbf{pandas}
            \item Charts with \textbf{matplotlib} (and a little seaborn)
            \item Feature types: numerical vs categorical
            \item Train / test thinking; simple metrics
            \item Models with \textbf{scikit-learn}: k-NN, linear \& logistic regression, trees
            \item A mini-project and a first look at ethics / limits of AI
        \end{itemize}
    \end{frame}

    \begin{frame}{What we will \emph{not} deep-dive}
        \begin{itemize}
            \item Heavy calculus or ``deriving gradients''
            \item Matrix theory and advanced linear algebra
            \item Building ChatGPT-scale models from scratch
            \item Transformer internals and research papers
        \end{itemize}
        \vspace{0.5em}
        \commonmistake{Thinking you must be a maths olympiad winner to start.
        Curiosity + careful practice is enough for this course.}
    \end{frame}

    \begin{frame}{How success looks after a few weeks}
        You will be able to:
        \begin{enumerate}
            \item Open a CSV, explore it, and plot something meaningful
            \item Explain features and a target in a supervised problem
            \item Train a simple model and talk about whether it seems useful
        \end{enumerate}
        \vspace{0.5em}
        Fancy architectures can wait for a later course.
    \end{frame}

%--- Quiz 2 ---
    \begin{frame}{Quiz 2 --- Course scope (true/false)}
        \quizbox{True or false?\\[0.4em]
        In this course we will derive neural-network training maths (gradients) from scratch.}
        \begin{enumerate}
            \item True
            \item False
        \end{enumerate}
        \quizpause
    \end{frame}

    \begin{frame}{Answer --- Quiz 2}
        \answerbox{\textbf{False.} We stay practical: libraries, intuition, and clear experiments---not heavy derivations.}
    \end{frame}

%======================================================================
    \section{Supervised learning}
%======================================================================

    \begin{frame}{The core idea}
        \begin{block}{Supervised learning}
            We have \textbf{examples} with known \textbf{answers} (labels).
            We want a model that predicts the answer for \textbf{new} examples.
        \end{block}
        \vspace{0.6em}
        \begin{itemize}
            \item \textbf{Features / inputs:} what we know about each example
            \item \textbf{Label / target:} the answer we want to predict
            \item \textbf{Model:} the learned rule (often a function from features $\rightarrow$ prediction)
        \end{itemize}
    \end{frame}

    \begin{frame}{Two common flavours}
        \begin{columns}[T]
            \begin{column}{0.48\textwidth}
                \begin{block}{Classification}
                    Predict a \textbf{category}\\[0.4em]
                    {\small spam / not spam\\
                    pass / fail\\
                    cat / dog / bird}
                \end{block}
            \end{column}
            \begin{column}{0.48\textwidth}
                \begin{block}{Regression}
                    Predict a \textbf{number}\\[0.4em]
                    {\small house price\\
                    temperature\\
                    exam score}
                \end{block}
            \end{column}
        \end{columns}
        \vspace{0.6em}
        Same big idea; different kind of answer.
    \end{frame}

    \begin{frame}{A tiny table (imagine a spreadsheet)}
        \begin{center}
            \begin{tabular}{@{}cccc@{}}
                \toprule
                \textbf{Size (m$^2$)} & \textbf{Rooms} & \textbf{City centre?} & \textbf{Price (k)} \\
                \midrule
                50 & 2 & yes & 180 \\
                80 & 3 & no  & 210 \\
                65 & 2 & yes & 200 \\
                \ldots & \ldots & \ldots & \ldots \\
                \bottomrule
            \end{tabular}
        \end{center}
        \vspace{0.6em}
        \begin{itemize}
            \item Features: size, rooms, city centre
            \item Label: price
            \item Task type: \textbf{regression} (a number)
        \end{itemize}
    \end{frame}

    \begin{frame}{Training vs using the model}
        \begin{enumerate}
            \item \textbf{Train:} show the model many examples with correct labels
            \item \textbf{Predict:} give new features; get a predicted label
            \item \textbf{Check:} compare predictions to reality on data the model did \emph{not} train on
        \end{enumerate}
        \vspace{0.6em}
        We will practise the ``don't cheat with the test data'' idea in a later week.
    \end{frame}

    \begin{frame}{Not everything is supervised}
        Quick awareness (details later if at all):
        \begin{itemize}
            \item \textbf{Unsupervised:} find structure without labels (e.g.\ group similar customers)
            \item \textbf{Generative AI:} create text or images (powerful tools; different workflow)
        \end{itemize}
        \vspace{0.5em}
        This semester's backbone is \textbf{supervised learning on tables}.
    \end{frame}

%--- Quiz 3 ---
    \begin{frame}{Quiz 3 --- Supervised learning}
        \quizbox{You predict whether an email is spam or not spam from its words.
        What kind of supervised task is this?}
        \begin{enumerate}
            \item Regression
            \item Classification
            \item Neither---there are no features
            \item Unsupervised clustering only
        \end{enumerate}
        \quizpause
    \end{frame}

    \begin{frame}{Answer --- Quiz 3}
        \answerbox{\textbf{2.} Classification --- the label is a category (spam / not spam).}
        \vspace{0.4em}
        If we predicted ``how many spam emails arrive tomorrow'' (a count), that would lean toward regression.
    \end{frame}

    \begin{frame}{Quiz 3b --- Features and labels}
        \quizbox{In the housing table, \textbf{price} is best described as:}
        \begin{enumerate}
            \item A feature used as input
            \item The label / target we want to predict
            \item The programming language
            \item A type of chart
        \end{enumerate}
        \quizpause
    \end{frame}

    \begin{frame}{Answer --- Quiz 3b}
        \answerbox{\textbf{2.} The label / target --- what we try to predict from the other columns.}
    \end{frame}

%======================================================================
    \section{Tools}
%======================================================================

    \begin{frame}{Our toolkit this semester}
        \begin{itemize}
            \item \textbf{Python} --- the language
            \item \textbf{Jupyter} (or Google Colab) --- notebooks: text + code + results together
            \item \textbf{pandas} --- tables (DataFrames)
            \item \textbf{matplotlib} --- plots
            \item \textbf{scikit-learn} --- classical ML models with a consistent API
        \end{itemize}
        \vspace{0.5em}
        Together, these are the ``framework'' feeling: shared patterns like \pycode{fit} / \pycode{predict}.
    \end{frame}

    \begin{frame}{Why Jupyter?}
        \begin{itemize}
            \item Run a small piece of code; see output immediately
            \item Mix explanations and results (great for learning and projects)
            \item Easy to share a \pycode{.ipynb} file or use Colab in the browser
        \end{itemize}
        \vspace{0.5em}
        \trythis{In practice today: create a notebook, run a cell with
        \pycode{print("Hello, AI Frameworks!")}.}
    \end{frame}

    \begin{frame}[fragile]{A first taste of Python in a notebook}
        \begin{lstlisting}
print("Hello, AI Frameworks!")

# A tiny prediction "by hand" (no ML library yet)
size = 70          # m^2
guess_price = 2.5 * size + 40   # made-up rule
print("Guessed price (k):", guess_price)
        \end{lstlisting}
        \vspace{0.3em}
        Later, libraries will learn better rules from real data---same idea, less guesswork.
    \end{frame}

    \begin{frame}{Libraries = don't reinvent everything}
        \begin{itemize}
            \item Professionals rarely write sorting, plotting, or k-NN from scratch in coursework
            \item Libraries give tested tools and a common language with teammates
            \item Your job: ask good questions, prepare data, interpret results
        \end{itemize}
        \vspace{0.4em}
        \commonmistake{Thinking ``using a library means I'm not really learning.''
        Using tools well \emph{is} the skill---plus understanding what they do.}
    \end{frame}

%--- Quiz 4 ---
    \begin{frame}{Quiz 4 --- Tools}
        \quizbox{Which tool is mainly for working with \textbf{tables} of data in Python?}
        \begin{enumerate}
            \item matplotlib
            \item pandas
            \item HTML
            \item Photoshop
        \end{enumerate}
        \quizpause
    \end{frame}

    \begin{frame}{Answer --- Quiz 4}
        \answerbox{\textbf{2.} pandas --- DataFrames for tabular data.\\[0.3em]
        matplotlib is for plots; we will use both often.}
    \end{frame}

%======================================================================
    \section{Wrap-up}
%======================================================================

    \begin{frame}{Key takeaways}
        \begin{enumerate}
            \item AI is broad; ML learns from data; data science includes exploring data
            \item This course is practical and light on formal maths
            \item Supervised learning: features + labels $\rightarrow$ predictions on new cases
            \item Classification predicts categories; regression predicts numbers
            \item Python + Jupyter + pandas / matplotlib / scikit-learn are our tools
        \end{enumerate}
    \end{frame}

    \begin{frame}{Next up}
        \begin{block}{Practice (today)}
            Environment check · first notebook · predict house price by hand
        \end{block}
        \vspace{0.5em}
        \begin{block}{Next lecture (Week 2)}
            Python enough for data work: lists, dicts, loops, functions
        \end{block}
    \end{frame}

    \begin{frame}{Questions?}
        \begin{center}
            {\Large What is still fuzzy: AI vs ML, labels, or the tools?}
        \end{center}
    \end{frame}

\end{document}